\begin{thebibliography}{99}
\bibitem{ABHM19}M. Akman, S. Bortz, S. Hofmann \& J. M. Martell - “Rectifiability, interior approximation and harmonic measure”, Ark. Mat. 57 (2019) no. 1, p. 1-22 | DOI | MR | Zbl
\bibitem{AHMMMTV2016GAFA}J. Azzam, S. Hofmann, J. M. Martell, S. Mayboroda, M. Mourgoglou, X. Tolsa \& A. L. Volberg - “Rectifiability of harmonic measure”, Geom. Funct. Anal. 26 (2016) no. 3, p. 703-728 | DOI | Zbl | MR
\bibitem{AHMMT2020Inventiones}J. Azzam, S. Hofmann, J. M. Martell, M. Mourgoglou \& X. Tolsa - “Harmonic measure and quantitative connectivity: geometric characterization of the $L^p$-solvability of the Dirichlet problem”, Invent. Math. 222 (2020) no. 3, p. 881-993 | DOI | Zbl | MR
\bibitem{Aikawa01}H. Aikawa - “Boundary Harnack principle and Martin boundary for a uniform domain”, J. Math. Soc. Japan 53 (2001) no. 1, p. 119-145 | DOI | MR | Zbl
\bibitem{Ancona1986Strong}A. Ancona - “On strong barriers and an inequality of Hardy for domains in ${\mathbf{R}}^n$”, J. London Math. Soc. (2) 34 (1986) no. 2, p. 274-290 | DOI | MR | Zbl
\bibitem{Azzam2020}J. Azzam - “Dimension drop for harmonic measure on Ahlfors regular boundaries”, Potential Anal. 53 (2020) no. 3, p. 1025-1041 | DOI | MR | Zbl
\bibitem{Azzam2021}J. Azzam - “Semi-uniform domains and the $A_\infty $ property for harmonic measure”, Internat. Math. Res. Notices (2021) no. 9, p. 6717-6771 | DOI | MR | Zbl
\bibitem{Batakis1996}A. Batakis - “Harmonic measure of some Cantor type sets”, Ann. Acad. Sci. Fenn. Math. 21 (1996) no. 2, p. 255-270 | MR | Zbl
\bibitem{BatakisZdunik2015}A. Batakis \& A. Zdunik - “Hausdorff and harmonic measures on non-homogeneous Cantor sets”, Ann. Acad. Sci. Fenn. Math. 40 (2015) no. 1, p. 279-303 | DOI | MR | Zbl
\bibitem{Carleson1985}L. Carleson - “On the support of harmonic measure for sets of Cantor type”, Ann. Acad. Sci. Fenn. Ser. A I Math. 10 (1985), p. 113-123 | DOI | MR | Zbl
\bibitem{Dah}B. E. J. Dahlberg - “Estimates of harmonic measure”, Arch. Rational Mech. Anal. 65 (1977) no. 3, p. 275-288 | DOI | MR | Zbl
\bibitem{DEM21}G. David, M. Engelstein \& S. Mayboroda - “Square functions, nontangential limits, and harmonic measure in codimension larger than 1”, Duke Math. J. 170 (2021) no. 3, p. 455-501 | DOI | MR | Zbl
\bibitem{DFM-Memoirs}G. David, J. Feneuil \& S. Mayboroda - Elliptic theory for sets with higher co-dimensional boundaries, Mem. Amer. Math. Soc., vol. 274, no. 1346, American Mathematical Society, Providence, RI, 2021 | DOI
\bibitem{DFM-Mixed}G. David, J. Feneuil \& S. Mayboroda - Elliptic theory in domains with boundaries of mixed dimension, Astérisque, vol. 442, Société Mathématique de France, Paris, 2023
\bibitem{DJ}G. David \& D. Jerison - “Lipschitz approximation to hypersurfaces, harmonic measure, and singular integrals”, Indiana Univ. Math. J. 39 (1990) no. 3, p. 831-845 | DOI | MR | Zbl
\bibitem{DavidMayboroda-good}G. David \& S. Mayboroda - “Good elliptic operators on Cantor sets”, Adv. Math. 383 (2021), article ID 107687, 21 pages | DOI | MR | Zbl
\bibitem{DM-Compactness}G. David \& S. Mayboroda - “Approximation of Green functions and domains with uniformly rectifiable boundaries of all dimensions”, Adv. Math. 410 (2022), article ID 108717, 52 pages | DOI | MR | Zbl
\bibitem{GM05}J. B. Garnett \& D. E. Marshall - Harmonic measure, New Math. Monogr., vol. 2, Cambridge University Press, Cambridge, 2005 | DOI
\bibitem{JonesWolff}P. W. Jones \& T. H. Wolff - “Hausdorff dimension of harmonic measures in the plane”, Acta Math. 161 (1988) no. 1-2, p. 131-144 | DOI | MR | Zbl
\bibitem{MakarovVolberg1986}N. Makarov \& A. L. Volberg - “On the harmonic measure of discontinuous fractals”, 1986, LOMI preprints
\bibitem{Tolsa2023ADreg}X. Tolsa - “The dimension of harmonic measure on some AD-regular flat sets of fractional dimension”, 2023 | arXiv
\bibitem{Volberg1992}A. L. Volberg - “On the harmonic measure of self-similar sets on the plane”, in Harmonic analysis and discrete potential theory (Frascati, 1991), Plenum, New York, 1992, p. 267-280 | DOI | MR
\bibitem{Volberg1993}A. L. Volberg - “On the dimension of harmonic measure of Cantor repellers”, Michigan Math. J. 40 (1993) no. 2, p. 239-258 | DOI | MR | Zbl
\bibitem{Volberg2022onephase}A. L. Volberg - “One phase problem for two positive harmonic functions: below the codimension 1 threshold”, 2022 | arXiv
\bibitem{Zdunik1997}A. Zdunik - “Harmonic measure on the Julia set for polynomial-like maps”, Invent. Math. 128 (1997) no. 2, p. 303-327 | DOI | MR | Zbl
\end{thebibliography}
